\apphead{Appendix 2}{Frobenius}

In this appendix we will explain the relation between the Frobenius\index{Frobenius}
endomorphism on p-adic modular forms\index{p-adic modular form} and the action of Frobenius on the de Rham
cohomology of "the" universal elliptic curve\index{elliptic curve}.

\kaitem{A2.0} Let $R$ be a p-adically complete ring, $E/R$ an elliptic curve
which modulo $p$ has invertible Hasse invariant\index{Hasse invariant}, and $H\subset E$ its canonical
subgroup. Let $E'=E/H$, and let $\pi\colon E\to E'$ denote the projection. Then
$\pi$ induces an R-morphism
$\pi^*\colon H^1_{\dR}(E'/R)\to H^1_{\dR}(E/R)$. Suppose now that
$R=M(W(\F_q),1,n,0)$, the ring of p-adic modular functions of level $n$ defined
over $W(\F_q)$, where $q$ is a power of $p$ such that $q\equiv 1 \bmod n$. Let
$E/R$ be the universal curve with level $n$ structure, such that Hasse is
invertible mod $p$. As $E'=E/H$ is a curve over $R$ with level $n$ structure and
Hasse invertible mod $p$, it is "classified" by a unique homomorphism
$\varphi\colon R\to R$ such that $E'=E^{(\varphi)}$. This homomorphism $\varphi$
is precisely the Frobenius endomorphism of the ring $M(W(\F_q),1,n,0)$ defined
in \cite{dwork69} (the "Deligne-Tate mapping"). The induced homomorphism
$\pi^*\colon H^1_{\dR}(E'/R)=H^1_{\dR}(E^{(\varphi)}/R)=(H^1(E/R))^{(\varphi)}
\to H^1_{\dR}(E/R)$ gives a $\varphi$-linear endomorphism of $H^1_{\dR}(E/R)$,
which we denote $F(\varphi)=\pi^*\circ\varphi^{-1}$ (to be compatible with the
notations of \cite{katz73}). Because $\pi^*$ is induced by an R-morphism $E\to E'$, the
endomorphism $F(\varphi)$ \uline{respects} the Hodge filtration
\[
0\to\w\to H^1_{\dR}(E/R)\to\w^{-1}\to 0 ,
\]
and thus induces $\varphi$-linear endomorphisms (still noted $F(\varphi)$) of
$\w$ and of $\w^{-1}$.

\secmark{A2.1}{Relation of the De Rham and p-adic modular Frobenii}

\begin{kathm}{Lemma}{(A2.1).}
On $\w$, $F(\varphi)=p\varphi$; on $\w^{-1}$, $F(\varphi)=\varphi$.
\end{kathm}

\begin{kaproof}
(We will suppress the level n structures, for simplicity.) Let $f$ be a section
of $\w$. Then $f(E,\omega)\cdot\omega$ is a section of $\Omega^1_{E/R}$. By
definition, $\varphi(f)$ is the section $f(E/H,\pi^*(\omega))\cdot\omega$ of
$\Omega^1_{E/R}$. Because $\pi$ is \'etale and $E/H=E^{(\varphi)}$, we have
$\pi^*(\omega)=\lambda\cdot\omega^{(\varphi)}$ with $\lambda$ invertible in $R$.
Thus $f(E/H,\pi^*(\omega))\cdot\omega=f(E^{(\varphi)},\lambda\omega^{(\varphi)})
\cdot\omega=\lambda^{-1}\cdot\varphi(f(E,\omega))\cdot\omega$. On the other hand,
\[
F(\varphi)(f(E,\omega))\xlongequal{\text{defn}}
\pi^*((f(E,\omega)\cdot\omega)^{(\varphi)})
=\varphi(f(E,\omega))\cdot\pi^*(\omega^{(\varphi)})
=\varphi(f(E,\omega))\cdot\frac{p\omega}{\lambda},
\]
[the last equality because
$p\omega=[p]^*(\omega)=\pi^*(\widetilde{\pi}^*(\omega))
=\pi^*(\lambda\cdot\omega^{(\varphi)})=\lambda\cdot\pi^*(\omega^{(\varphi)})$].
Thus $F(\varphi)=p\varphi$ as $\varphi$-linear endomorphism.

Similarly, for $\w^{-1}$, a section $f$ is a section
$f(E,\omega)\cdot\omega^{-1}$ of $H^1(E,\cO_E)$, and $\varphi(f)$ is the section
$f(E/H,\pi^*(\omega))\cdot\omega^{-1}$ of $H^1(E,\cO_E)$. But as before
$E/H=E^{(\varphi)}$, $\pi^*(\omega)=\lambda\omega$ with $\lambda$ invertible in
$R$, and so $\varphi(f)$ is the section $\lambda\cdot\varphi(f(E,\omega))\cdot
\omega^{-1}$. But
$F(\varphi)(f(E,\omega)\cdot\omega^{-1})
=\pi^*(\varphi(f(E,\omega))\cdot(\omega^{-1})^{(\varphi)})
=\varphi(f(E,\omega))\cdot\pi^*((\omega^{-1})^{(\varphi)})$. So we must show
that $\pi^*((\omega^{-1})^{(\varphi)})=\lambda\cdot\omega^{-1}$, or by Serre
duality, that $\pi^*(\omega^{(\varphi)})=\lambda\cdot\omega$, which was the
\uline{definition} of $\lambda$.
\end{kaproof}

\sechead{A2.2}{Calculation at $\infty$}

The canonical subgroup\index{canonical subgroup} of $\mathrm{Tate}(q)$ over $\Z((q))$ is $\mu_p$, and the
quotient is $\mathrm{Tate}(q^p)=\mathrm{Tate}(q)^{(\varphi)}$, where
$(qf)(q)=f(q^p)$. Thus we also have a $\varphi$-linear endomorphism of
$H^1_{\dR}(\mathrm{Tate}(q)/\Z((q)))$. Passing to $\C((q))$ and viewing the
situation analytically, $\omega_{\mathrm{can}}$ becomes the differential
$2\pi i\,\dd z$ on $\C/\Z+\Z\tau$, and the canonical subgroup becomes
$\frac{1}{p}\Z/\Z$. The quotient is
$\C/\frac{1}{p}\Z+\Z\tau\xrightarrow{\varphi}\C/\Z+\Z\cdot p\tau$. In terms of
the bases $\gamma_i(\tau)$, $i=1,2$ and $\gamma_i(p\tau)$, $i=1,2$ of $H_1$, we
have $\pi\gamma_1(\tau)=\gamma_1(p\tau)$, $\pi(\gamma_2(\tau))=p\gamma_2(p\tau)$.
It follows that
$\pi^*(\omega_{\mathrm{can}}(q^p))=\pi^*((\omega_{\mathrm{can}}(q))^{(\varphi)})
=p\cdot\omega_{\mathrm{can}}(q)$ because both have the same periods:

\begin{kadispb}{A2.2.1}
\begin{cases}
\displaystyle\int_{\gamma_1(\tau)}\pi^*(\omega_{\mathrm{can}}(q^p))
=\int_{\pi\gamma_1}\omega_{\mathrm{can}}(q^p)
=\int_{\gamma_1(p\tau)}\omega_{\mathrm{can}}(q^p)=p\tau \\[6pt]
\displaystyle\int_{\gamma_2(\tau)}\pi^*(\omega_{\mathrm{can}}(q^p))
=\int_{\pi\gamma_2}\omega_{\mathrm{can}}(q^p)
=\int_{p\gamma_2(p\tau)}\omega_{\mathrm{can}}(q^p)=p .
\end{cases}
\end{kadispb}

By functionality, $F(\varphi)$ respects the Gauss-Manin connection\index{Gauss--Manin connection}, and
$\nabla(\theta)(\omega_{\mathrm{can}})=\frac{-P}{12}\omega_{\mathrm{can}}
+\eta_{\mathrm{can}}$ is the \uline{unique} (up to scalars) element of
$H^1_{\dR}(\mathrm{Tate}(q)/\Z((q)))$ killed by $\nabla(\theta)$ (as a direct
calculation shows --- indeed by \eqref{ka:A1.4.6}, this rank two differential equation
over $\C$ has non-trivial (unipotent) monodromy around $q=0$, hence has at most
\uline{one} solution which is single-valued at $q=0$). It follows that

\begin{kadisp}{A2.2.2}
F(\varphi)(\nabla(\theta)(\omega_{\mathrm{can}}))=a\cdot\nabla(\theta)\omega_{\mathrm{can}}
\end{kadisp}

for some $a\in\Z$; explicitly,

\begin{kadisp}{A2.2.3}
\frac{-\varphi(P)}{12}\pi^*(\omega_{\mathrm{can}}^{(\varphi)})
+\pi^*(\eta_{\mathrm{can}}^{(\varphi)})
=\frac{-aP}{12}\omega_{\mathrm{can}}+a\cdot\eta_{\mathrm{can}} ,
\end{kadisp}

whence

\begin{kadisp}{A2.2.4}
F(\varphi)(\eta_{\mathrm{can}})
=\frac{P\cdot\varphi(P)-aP}{12}\omega_{\mathrm{can}}+a\cdot\eta_{\mathrm{can}} .
\end{kadisp}

Because $\omega_{\mathrm{can}}$ and $\nabla(\theta)\omega_{\mathrm{can}}$ give a
\uline{base} of $H^1$ such that
$\omega_{\mathrm{can}}\wedge\nabla(\theta)\omega_{\mathrm{can}}$ is a
\uline{constant} base of $H^2$, the fact that $\pi$ has degree $p$ shows that
$a=1$, so

\begin{kadisp}{A2.2.5}
F(\varphi)(\eta_{\mathrm{can}})
=\frac{P\cdot\varphi(P)-P}{12}\omega_{\mathrm{can}}+\eta_{\mathrm{can}} .
\end{kadisp}

Thus the \uline{matrix} of $F(\varphi)$ on
$H^1(\mathrm{Tate}(q)/\Z[1/6]((q)))$ is given by

\begin{kadisp}{A2.2.6}
\begin{pmatrix}F(\varphi)(\omega_{\mathrm{can}})\\ F(\varphi)(\eta_{\mathrm{can}})\end{pmatrix}
=\begin{pmatrix} p & 0 \\ \frac{p\cdot\varphi(P)-P}{12} & 1\end{pmatrix}
\begin{pmatrix}\omega_{\mathrm{can}}\\ \eta_{\mathrm{can}}\end{pmatrix} .
\end{kadisp}

To give formulas valid over $\Z((q))$, we use the base
$\omega_{\mathrm{can}},\nabla(\theta)(\omega_{\mathrm{can}})$ of
$H^1_{\dR}(\mathrm{Tate}(q)/\Z((q)))$; we have

\begin{kadisp}{A2.2.7}
\begin{pmatrix}F(\varphi)(\omega_{\mathrm{can}})\\ F(\varphi)(\nabla(\theta)(\omega_{\mathrm{can}}))\end{pmatrix}
=\begin{pmatrix} p & 0 \\ 0 & 1\end{pmatrix}
\begin{pmatrix}\omega_{\mathrm{can}}\\ \nabla(\theta)(\omega_{\mathrm{can}})\end{pmatrix} .
\end{kadisp}

\sechead{A2.3}{The "canonical direction" in $H^1_{\dR}$ (a special case of \cite{katz73}, \cite{dwork71b})}

We return to the universal situation $R=M(W(\F_q),1,n,0)$, $E/R$
uni\-versal. In terms of a base $\omega,\eta$ of $H^1_{\dR}(E/R)$ adopted to the
Hodge filtration, the matrix of $F(\varphi)$ has the shape:

\begin{kadisp}{A2.3.1}
F(\varphi)\begin{pmatrix}\omega\\ \eta\end{pmatrix}
=\begin{pmatrix} p/\lambda & 0 \\ C & \lambda\end{pmatrix}
\begin{pmatrix}\omega\\ \eta\end{pmatrix}
\qquad\text{with}\quad
\begin{cases}\lambda\in R\ \ \uline{invertible}\\ C\in R\end{cases} .
\end{kadisp}

An argument of successive approximation shows that there is a unique element
$f\in R$ such that $F(\varphi)(\eta+f\omega)\in R\cdot(\eta+f\omega)$; indeed

\begin{kadisp}{A2.3.2}
F(\varphi)(\eta+f\omega)=c\omega+\lambda\eta+\varphi(f)\frac{p}{\lambda}\cdot\omega
=\Big\{c+\frac{p}{\lambda}\varphi(f)\Big\}\omega+\lambda\eta ,
\end{kadisp}

so we want $f\in R$ to satisfy

\begin{kadisp}{A2.3.3}
\begin{cases}
c+\dfrac{p}{\lambda}\varphi(f)=\lambda f \\[6pt]
\text{i.e.}\quad f=\dfrac{c}{\lambda}+\dfrac{p}{\lambda^2}\varphi(f) .
\end{cases}
\end{kadisp}

Let us define a mapping $T\colon R\to R$ by
$T(f)=\frac{c}{\lambda}+\frac{p}{\lambda^2}\varphi(f)$. It is immediate that $T$
is a \uline{contraction mapping} of $R$ in its p-adic topology, so has a unique
fixed point $\lim T^n(0)$, which is explicitly given by

\begin{kadisp}{A2.3.4}
f=\frac{c}{\lambda}+\sum_{n\geq 1}
\frac{p^n\,\varphi^{\frac{n(n-1)}{2}}(1/\lambda)\cdot\varphi^n(c)}
{\varphi^{\frac{n(n+1)}{2}}(\lambda)} .
\end{kadisp}

Of course, the choice of base is not canonical, nor need there exist a
\uline{global} basis (over all of $R$), but the given construction does
construct an $F(\varphi)$-\uline{splitting} of the Hodge filtration

\begin{equation}\tag{A2.3.5}\label{ka:A2.3.5}
\begin{tikzcd}[row sep=tiny, column sep=large, ampersand replacement=\&]
0 \arrow[r] \& \w \arrow[r] \& H^1_{\dR}(E/R) \arrow[r]
  \arrow[ll, dashed, bend right=40, looseness=1.4] \& \w^{-1} \arrow[r] \& 0 .
\end{tikzcd}
\end{equation}

Looking at $\infty$, we see that in terms of the base
$\omega_{\mathrm{can}},\allowbreak\nabla(\theta)(\omega_{\mathrm{can}})$ of
$H^1_{\dR}(\mathrm{Tate}(q)/\Z((q)))$, we have simply "constructed" the vector
$\nabla(\theta)(\omega_{\mathrm{can}})$, which is indeed fixed by $F(\varphi)$.
Hence we have proven

\begin{kathm}{Theorem}{A2.3.6.}
(Dwork) Let $\Mbar_n(W(\F_q),1)$ denote the formal scheme over $W(\F_q)$ which
mod $p^m$ is the open subset of $\Mbar_n\tens W_m(\F_q)$ where $E_{p-1}$ is
invertible. The locally free rank two module on $M_n(W(\F_q),1)$ given by
$H^1_{\dR}(E/M_n(W(\F_q),1))$ admits a locally free extension
$H^1_{\dR}(E/\Mbar_n(W(\F_q),1))$ which along any cusp is the
$W(\F_q)[[q]]$ submodule of $H^1(\mathrm{Tate}(q^n)/W(\F_q)((q)))$ spanned by
$\omega_{\mathrm{can}}$ and by $\nabla(\theta)(\omega_{\mathrm{can}})$. The
Gauss-Manin connection\index{Gauss--Manin connection} over $M_n(W(\F_q),1)$ extends to a "connection with
logarithmic poles" over $\Mbar_n(W(\F_q),1)$, and the $\varphi$-linear
endomorphism $F(\varphi)$ over $M_n(W(\F_q),1)$ extends to a $\varphi$-linear
endomorphism, still noted $F(\varphi)$, over all of $\Mbar_n(W(\F_q),1)$
(cf.\eqref{ka:A2.2.6} and \eqref{ka:A2.2.7} for the explicit formulas \uline{defining} these
extensions). There is a canonical $F(\varphi)$-stable splitting of the Hodge
filtration
\[
0\longrightarrow\w\longrightarrow H^1_{\dR}(E/\Mbar_n(W(\F_q),1))
\longrightarrow\w^{-1}\longrightarrow 0 ,
\]
(the \uline{image} of which we denote $U\subset H^1_{\dR}(E/\Mbar_n(W(\F_q),1))$;
it is a horizontal (by unicity!) $F(\varphi)$-stable rank one submodule).
\end{kathm}

\sechead{A2.4}{P as a p-adic modular form of weight 2}

Suppose now that $p\neq 2,3$. Let $R$ be any ring in which $p$ is
nil\-potent, $E/R$ an elliptic curve\index{elliptic curve} whose Hasse invariant\index{Hasse invariant} modulo $p$ is
invertible, and $U\subset H^1_{\dR}(E/R)$ the inverse image of the canonical
rank one submodule constructed above. (Strictly speaking, we must first choose a
level n structure for some $n\geq 3$ prime to $p$ defined over an \'etale
over-ring $R'$ of $R$, and check that the $U$ obtained in
$H^1(E_{R'}/R')$ descends to a $U\subset H^1(E/R)$ which is independent of
choices.) Let $\omega$ be a nowhere-vanishing differential on $E/R$ (which in
any case exists locally on $R$), and let $\eta$ be the corresponding
differential of the second kind (i.e.\ $\omega=\frac{\dd X}{Y}$,
$\eta=\frac{X\,\dd X}{Y}$ as explained in \eqref{ka:A1.2.4}). Because
$H^1=R\cdot\omega+U$, we see that if $u\in U$ is a base of $U$ (which in any
case exists locally on $R$) then the de Rham cup-product
$\langle\omega,u\rangle$ is \uline{invertible} on $R$. We may then define a
"function" $\widetilde{P}$ by the formula

\begin{kadisp}{A2.4.1}
\widetilde{P}(E/R,\omega)=12\underbrace{\frac{\langle\eta,u\rangle}{\langle\omega,u\rangle}}_{\text{for any base } u \text{ of } U} .
\end{kadisp}

Clearly the right-hand expression is independent of the choice of base $u$ of
$U$, and the effect of replacing $\omega$ by $\lambda\omega$, $\lambda\in R^\times$
is to replace $\eta$ by $\lambda^{-1}\eta$, hence
$\widetilde{P}(E/R,\lambda\omega)=\lambda^{-2}\widetilde{P}(E/R,\omega)$. Hence
$\widetilde{P}$ is a p-adic modular form\index{p-adic modular form} of weight two and level one. Its
q-expansion is

\begin{kadispb}{A2.4.2}
\widetilde{P}(\mathrm{Tate}(q),\omega_{\mathrm{can}})
=12\frac{\langle\eta_{\mathrm{can}},u\rangle}{\langle\omega_{\mathrm{can}},u\rangle}
=12\frac{\langle\eta_{\mathrm{can}},\nabla(\theta)(\omega_{\mathrm{can}})\rangle}
{\langle\omega_{\mathrm{can}},\nabla(\theta)(\omega_{\mathrm{can}})\rangle}
\end{kadispb}

because, formally at $\infty$, $U$ is spanned by
$\nabla(\theta)(\omega_{\mathrm{can}})$. If we denote by $P(q)$ the \uline{series}
$1-24\sum_{n\geq 1}\sigma_1(n)q^n$, then by \eqref{ka:A1.3.16} we have
$\nabla(\theta)(\omega_{\mathrm{can}})=\frac{-P(q)}{12}\omega_{\mathrm{can}}
+\eta_{\mathrm{can}}$. Substituting into \eqref{ka:A2.4.2} gives

\begin{kadisp}{A2.4.3}
\widetilde{P}(\mathrm{Tate}(q),\omega_{\mathrm{can}})
=12\frac{\big\langle\eta_{\mathrm{can}},\frac{-P(q)}{12}\omega_{\mathrm{can}}+\eta_{\mathrm{can}}\big\rangle}
{\big\langle\omega_{\mathrm{can}},\frac{-P(q)}{12}\omega_{\mathrm{can}}+\eta_{\mathrm{can}}\big\rangle}
=\frac{12\cdot P(q)}{12}\frac{\langle\eta_{\mathrm{can}},-\omega_{\mathrm{can}}\rangle}
{\langle\omega_{\mathrm{can}},\eta_{\mathrm{can}}\rangle}
=P(q) .
\end{kadisp}

This provides a modular proof that $P$ is p-adically modular.
